\subsection{ESTIMACIÓN DE DETERIORO Y LUCRO CESANTE.}


El valuador llevó a cabo la valuación del deterioro y lucro cesante marcario, consistente en:
\begin{enumerate}[A.]

\item \textbf{\textcolor{principal}{Estimación del valor unitario por producto de la inversión (marcaria) sujeta de deterioro}}.- Consiste en la inversión neta considerada como el costo de ventas aproximado para la adquisición de inventario(s) disponible(s) ligados a la marca Innominada 694670.
\item \textbf{\textcolor{principal}{Estimación del valor unitario del Lucro Cesante por producto}}.- Consistente en el flujo de caja libre que dejó de obtener la marca Burberry a través de su franquiciatario en México la sociedad “Horseferry México”, S.A. de C.V. (en lo sucesivo “HorseFerry”).

\end{enumerate}

La estimación financiera se llevó a cabo a partir de la información de negocios proporcionada por el solicitante, habiendo aplicado el marco teórico del Análisis de Presupuesto de Capital para calcular el FLUJO DE CAJA LIBRE (FCL) bajo la metodología conocida como ``Ebit pormenorizado sin depreciación'', atribuida al profesor Aswath Damodaran, así como el método de valor agregado estimado (EVA), atribuida al profesor Tomas Copeland; bajo el número 2 del Pentágono de McKinsey (Valor actual interno o ``\textit{value as is}''):

\begin{figure}[H]
\centering
\caption{Pent\'agono de Explotaci\'on de oportunidades\label{fig:hexagono}}
\includegraphics[width=9cm]{\rutaImagenes/pentagono}\\
Fuente:Valuation. Copeland Tom, Koller Tim y Murrin Jack.\\

John Wiley \& Sons. 2000.
\end{figure}
